%=====================================================================
\chapter{Artificial Neural Networks}\label{ch:ann}
%=====================================================================
\locator{3}{Part III --- Supervised Learning Algorithms}

\begin{outcomes}
By the end of this chapter you will be able to:
\begin{tight}
  \item \textbf{Compute} the output of a perceptron and \textbf{train} one with the
        perceptron rule.
  \item \textbf{Explain} linear separability and why a single perceptron cannot
        learn XOR.
  \item \textbf{Derive} the delta rule as gradient descent on squared error.
  \item \textbf{Perform} one forward and backward pass of backpropagation through a
        small multilayer network by hand.
  \item \textbf{Choose} activation functions and \textbf{apply} early stopping,
        weight decay and dropout against overfitting.
\end{tight}
\end{outcomes}

\begin{prereq}
\chref{math} (the chain rule, gradient descent), \chref{logreg} (the sigmoid, the
softmax, cross-entropy) and \chref{generalization} (overfitting, regularisation).
\end{prereq}

\begin{keyterms}
artificial neuron $\bullet$ perceptron $\bullet$ threshold unit $\bullet$
perceptron training rule $\bullet$ linear separability $\bullet$ XOR $\bullet$
delta rule $\bullet$ sigmoid unit $\bullet$ multilayer perceptron $\bullet$ hidden
layer $\bullet$ forward pass $\bullet$ backpropagation $\bullet$ epoch $\bullet$
activation function $\bullet$ ReLU $\bullet$ early stopping $\bullet$ weight decay
$\bullet$ dropout
\end{keyterms}

\section{From Neurons to Units}

A biological neuron receives signals from thousands of others, and fires if their
combined stimulation exceeds a threshold. Learning, roughly, changes the strength
of the connections. An \term{artificial neural network} keeps just that much of
the biology: many simple \emph{units}, each computing a weighted sum of its inputs
and passing it through a simple function, connected so that the outputs of some are
the inputs of others. Everything the network knows is in the connection weights,
and learning means adjusting them.

Networks are the right tool when the input is high-dimensional and raw --- in
project work, the full text of tickets and commit messages, or the complete event
history of every task in a tracker --- when the target function is complex and
unknown, when some noise is acceptable, when long training times are acceptable,
and when nobody needs to read the learned function. Networks of this kind, made
deep, are the whole of deep learning.

\section{The Perceptron}

\begin{definitionbox}[Perceptron]
A \term{perceptron} computes a weighted sum of its inputs and thresholds it:
\[
o(\mathbf{x}) =
\begin{cases} 1 & \text{if } w_0 + w_1 x_1 + \dots + w_d x_d > 0 \\ 0 & \text{otherwise} \end{cases}
\]
With a constant input $x_0 = 1$ the bias $w_0$ is just another weight, and the unit
computes $\operatorname{step}(\mathbf{w}\cdot\mathbf{x})$. Its decision boundary is the
hyperplane $\mathbf{w}\cdot\mathbf{x} = 0$ --- the same family as logistic regression
and the linear SVM.
\end{definitionbox}

\dsfig{%
\begin{tikzpicture}[font=\scriptsize,
  inp/.style={circle, draw=secondary, fill=secondary!10, minimum size=0.75cm, font=\scriptsize},
  a/.style={-{Stealth[length=2mm]}, gray!60, line width=0.9pt}]
\node[inp] (x0) at (0,1.8)  {$1$};
\node[inp] (x1) at (0,0.6)  {$x_1$};
\node[inp] (x2) at (0,-0.6) {$x_2$};
\node[font=\small] at (0,-1.45) {$\vdots$};
\node[inp] (xd) at (0,-2.2) {$x_d$};
\node[circle, draw=primary, line width=1.2pt, fill=primary!10, minimum size=1.3cm,
      font=\large] (sum) at (3.4,0) {$\sum$};
\node[rounded corners=3pt, draw=accent, line width=1.1pt, fill=accent!8, minimum width=1.5cm,
      minimum height=1.1cm] (act) at (6.0,0) {};
\draw[accent, line width=1.1pt] (5.45,-0.3) -- (6.0,-0.3) -- (6.0,0.3) -- (6.55,0.3);
\node[font=\tiny, accent, anchor=north] at (6.0,-0.6) {step (or sigmoid)};
\draw[a] (x0) -- node[above, font=\tiny, sloped] {$w_0$} (sum);
\draw[a] (x1) -- node[above, font=\tiny, sloped] {$w_1$} (sum);
\draw[a] (x2) -- node[above, font=\tiny, sloped] {$w_2$} (sum);
\draw[a] (xd) -- node[above, font=\tiny, sloped] {$w_d$} (sum);
\draw[a] (sum) -- node[above, font=\tiny] {$\mathbf{w}\cdot\mathbf{x}$} (act);
\draw[a] (act) -- (8.2,0) node[right, font=\scriptsize] {$o$};
\node[anchor=north west, align=left, text width=4.6cm, font=\scriptsize, text=gray!55!black]
     at (9.0,1.9)
     {\textbf{AND} with $w_0 = -0.8$, $w_1 = w_2 = 0.5$:\\[2pt]
      $(0,0) \to -0.8 \to 0$\\
      $(0,1) \to -0.3 \to 0$\\
      $(1,0) \to -0.3 \to 0$\\
      $(1,1) \to +0.2 \to 1$\\[4pt]
      Change $w_0$ to $-0.3$ and the same unit computes \textbf{OR}.};
\end{tikzpicture}}%
{A perceptron: a weighted sum of the inputs, then a threshold. The bias is a weight on a
constant input of 1. With the weights shown on the right it computes logical AND.}{perceptron}

\begin{definitionbox}[The Perceptron Training Rule]
Start from small (or zero) weights. For each training example with target $t$ and
current output $o$, update every weight:
\[
w_i \leftarrow w_i + \eta\,(t - o)\,x_i
\]
Correct outputs change nothing; a missed positive pushes the weights towards the
input, a false positive pushes them away. If the training examples are
\term{linearly separable} --- some hyperplane separates the classes perfectly ---
this procedure is guaranteed to find one in a finite number of steps, for any
sufficiently small learning rate $\eta$.
\end{definitionbox}

\begin{worked}[Training a Perceptron to Compute AND]
Inputs with bias $x_0 = 1$; weights $(w_0, w_1, w_2)$ start at $(0, 0, 0)$; $\eta = 0.5$.
The output is 1 only if $\mathbf{w}\cdot\mathbf{x} > 0$.

\smallskip
\textbf{Epoch 1.} $(0,0)$, $(0,1)$, $(1,0)$: net $= 0$, output 0, target 0 --- no change.
$(1,1)$: net $0$, output 0, target 1 --- update by $0.5 \times (1 - 0) \times (1, 1, 1)$:
$\mathbf{w} = (0.5,\ 0.5,\ 0.5)$.

\textbf{Epoch 2.} $(0,0)$: net $0.5$, output 1, target 0 --- subtract $0.5 \times (1,0,0)$:
$\mathbf{w} = (0,\ 0.5,\ 0.5)$. $(0,1)$: net $0.5$, output 1 --- subtract
$0.5 \times (1,0,1)$: $\mathbf{w} = (-0.5,\ 0.5,\ 0)$. $(1,0)$: net $0$, correct.
$(1,1)$: net $0$, output 0 --- add $0.5 \times (1,1,1)$: $\mathbf{w} = (0,\ 1,\ 0.5)$.

\textbf{Epochs 3--5} continue the same way, and at the end of epoch 5 the weights are
$\mathbf{w} = (-1,\ 1,\ 0.5)$. \textbf{Check:} $(0,0) \to -1$, $(0,1) \to -0.5$,
$(1,0) \to 0$, all output 0; $(1,1) \to +0.5$, output 1. Epoch 6 makes no changes:
converged. (Trace epochs 3 to 5 yourself: it is standard examination practice.)
\end{worked}

\section{Linear Separability and XOR}

A single perceptron draws one straight line. AND and OR can be separated by a line;
XOR --- true when exactly one input is true --- cannot: its two positive points
$(0,1)$ and $(1,0)$ sit on one diagonal and its two negatives on the other. No
perceptron, however trained, can represent XOR. Minsky and Papert's 1969 book made
this point so forcefully that neural network research stalled for over a decade.

\dsfig{%
\begin{tikzpicture}[font=\scriptsize,
  one/.style={circle, fill=secondary, inner sep=2.3pt},
  zero/.style={circle, draw=accent, line width=1.1pt, fill=white, inner sep=2.1pt}]
\foreach \ox/\ttl in {0/AND, 4.2/OR, 8.4/XOR}{
  \begin{scope}[shift={(\ox,0)}]
  \draw[gray!55, -{Stealth[length=1.6mm]}] (-0.3,0) -- (2.3,0) node[right, font=\tiny] {$x_1$};
  \draw[gray!55, -{Stealth[length=1.6mm]}] (0,-0.3) -- (0,2.3) node[above, font=\tiny] {$x_2$};
  \node[font=\scriptsize\bfseries, text=primary] at (1.0,-0.7) {\ttl};
  \end{scope}}
% AND
\node[zero] at (0,0) {}; \node[zero] at (0,1.6) {}; \node[zero] at (1.6,0) {}; \node[one] at (1.6,1.6) {};
\draw[greenhl!70!black, line width=1.2pt] (0.5,2.2) -- (2.2,0.5);
% OR
\node[zero] at (4.2,0) {}; \node[one] at (4.2,1.6) {}; \node[one] at (5.8,0) {}; \node[one] at (5.8,1.6) {};
\draw[greenhl!70!black, line width=1.2pt] (3.8,1.0) -- (5.2,-0.4);
% XOR
\node[zero] at (8.4,0) {}; \node[one] at (8.4,1.6) {}; \node[one] at (10.0,0) {}; \node[zero] at (10.0,1.6) {};
\draw[accent, line width=1.1pt, dashed] (8.0,0.9) -- (10.4,0.7);
\node[font=\tiny, accent, align=center] at (11.6,0.8) {no single line\\separates them};
\end{tikzpicture}}%
{AND and OR are linearly separable; XOR is not. A second layer of units, each drawing its
own line, is needed to represent it.}{xor}

The cure is a second layer. XOR is ``OR but not AND'': one hidden unit computes OR,
another computes NAND (not-AND), and an output unit computes AND of the two. Each
unit draws one line; together they carve out the diagonal. This is why networks have
hidden layers --- and the only problem is how to train them, since the perceptron
rule needs a target for every unit and the hidden units have none.

\section{Gradient Descent and the Delta Rule}

The perceptron rule cannot be extended to hidden units because the step function
has no useful derivative. The fix is to replace hard thresholds with smooth units
and train by gradient descent on an error measure.

\begin{definitionbox}[The Delta Rule]
For a \emph{linear} unit $o = \mathbf{w}\cdot\mathbf{x}$ and squared error
$E = \tfrac{1}{2}\sum_{d}(t_d - o_d)^2$ over the training examples, the gradient is
$\partial E/\partial w_i = -\sum_d (t_d - o_d)\,x_{id}$, so gradient descent updates
\[
w_i \leftarrow w_i + \eta\sum_d (t_d - o_d)\,x_{id}
\]
the \term{delta rule} (or LMS rule). It looks like the perceptron rule, but it is
gradient descent: it converges towards the minimum-error weights even when the data
is \emph{not} separable, which the perceptron rule does not. It is exactly the
gradient descent for linear regression of \chref{linreg}.
\end{definitionbox}

For a \term{sigmoid unit}, $o = \sigma(\mathbf{w}\cdot\mathbf{x})$, the chain rule adds
one factor, the derivative of the sigmoid, which has the convenient form
$\sigma'(z) = \sigma(z)\,(1 - \sigma(z)) = o(1 - o)$.

\section{Multilayer Networks and Backpropagation}

\begin{definitionbox}[Multilayer Perceptron]
A \term{multilayer perceptron} (MLP) arranges sigmoid (or other smooth) units in
layers: an input layer, one or more \term{hidden layers}, and an output layer. Each
unit takes the outputs of the previous layer, forms a weighted sum and applies its
activation. A network with one hidden layer of enough units can approximate any
continuous function to any accuracy; two hidden layers, any function at all.
\end{definitionbox}

\begin{definitionbox}[Backpropagation (for sigmoid units and squared error)]
For each training example:
\begin{tightnum}
  \item \textbf{Forward pass:} compute every unit's output $o$, layer by layer.
  \item \textbf{Output units:} $\delta_k = o_k(1 - o_k)(t_k - o_k)$.
  \item \textbf{Hidden units} (backwards, layer by layer):
        $\delta_h = o_h(1 - o_h)\sum_{k} w_{kh}\,\delta_k$ --- each hidden unit's share
        of the blame is the output deltas, weighted by its connections to them.
  \item \textbf{Update every weight} $w_{ji}$ from unit $i$ to unit $j$:
        $w_{ji} \leftarrow w_{ji} + \eta\,\delta_j\,x_{ji}$, where $x_{ji}$ is the input
        arriving along that connection.
\end{tightnum}
One sweep through the whole training set is an \term{epoch}.
\end{definitionbox}

\dsfig{%
\begin{tikzpicture}[font=\scriptsize,
  inp/.style={circle, draw=secondary, fill=secondary!10, minimum size=0.95cm},
  hid/.style={circle, draw=purplehl, fill=purplehl!10, minimum size=1.05cm, align=center},
  outu/.style={circle, draw=greenhl!70!black, fill=greenhl!12, minimum size=1.05cm, align=center},
  w/.style={font=\tiny, fill=white, inner sep=0.8pt},
  a/.style={-{Stealth[length=2mm]}, gray!60, line width=0.8pt}]
\node[inp] (x1) at (0,1.2) {$x_1 = 1$};
\node[inp] (x2) at (0,-1.2) {$x_2 = 0$};
\node[hid] (h1) at (4.2,1.4) {$h_1$\\[-1pt]{\tiny 0.622}};
\node[hid] (h2) at (4.2,-1.4) {$h_2$\\[-1pt]{\tiny 0.475}};
\node[outu] (o) at (8.2,0) {$o$\\[-1pt]{\tiny 0.658}};
\draw[a] (x1) -- node[w, pos=0.35] {0.4} (h1);
\draw[a] (x1) -- node[w, pos=0.25] {$-0.3$} (h2);
\draw[a] (x2) -- node[w, pos=0.25] {$-0.2$} (h1);
\draw[a] (x2) -- node[w, pos=0.35] {0.5} (h2);
\draw[a] (h1) -- node[w] {0.6} (o);
\draw[a] (h2) -- node[w] {0.8} (o);
\node[font=\tiny, purplehl, anchor=south] at (4.2,2.0) {bias 0.1};
\node[font=\tiny, purplehl, anchor=north] at (4.2,-2.0) {bias 0.2};
\node[font=\tiny, greenhl!45!black, anchor=south] at (8.2,0.6) {bias $-0.1$};
\draw[-{Stealth[length=2mm]}, accent, line width=1pt] (8.9,-0.3) -- (10.1,-0.3);
\node[font=\tiny, accent, anchor=west, align=left] at (10.2,-0.3) {target $t = 1$\\error $0.342$};
\draw[-{Stealth[length=2mm]}, accent!70, dashed, line width=1pt]
     (7.6,-0.9) .. controls (6.5,-2.4) and (5.2,-2.6) .. (4.7,-2.0);
\node[font=\tiny, accent!80!black, align=center] at (6.5,-2.75)
     {deltas flow backwards: $\delta_o = 0.0770$, $\delta_{h_1} = 0.0109$, $\delta_{h_2} = 0.0154$};
\end{tikzpicture}}%
{The network of the worked example, with the unit outputs of the forward pass written in
each unit. Errors computed at the output are passed backwards along the same connections
to assign blame to the hidden units.}{backprop-net}

\begin{worked}[One Step of Backpropagation]
The 2--2--1 network of \dsref{backprop-net}, all sigmoid units, input $(1, 0)$, target
$t = 1$, learning rate $\eta = 0.5$.

\smallskip
\textbf{Forward.}
\begin{tight}
  \item $h_1$: net $= 0.1 + 0.4(1) - 0.2(0) = 0.5$, output $\sigma(0.5) = 0.6225$.
  \item $h_2$: net $= 0.2 - 0.3(1) + 0.5(0) = -0.1$, output $\sigma(-0.1) = 0.4750$.
  \item $o$: net $= -0.1 + 0.6(0.6225) + 0.8(0.4750) = 0.6535$, output $\mathbf{0.6578}$.
        Error $t - o = 0.3422$.
\end{tight}
\textbf{Backward.}
\begin{tight}
  \item $\delta_o = 0.6578 \times (1 - 0.6578) \times 0.3422 = \mathbf{0.0770}$.
  \item $\delta_{h_1} = 0.6225 \times 0.3775 \times (0.6 \times 0.0770) = \mathbf{0.0109}$.
  \item $\delta_{h_2} = 0.4750 \times 0.5250 \times (0.8 \times 0.0770) = \mathbf{0.0154}$.
\end{tight}
\textbf{Update} ($\Delta w = \eta\,\delta\,x$).
\begin{tight}
  \item Output weights: from $h_1$: $0.6 + 0.5(0.0770)(0.6225) = 0.6240$; from $h_2$:
        $0.8 + 0.5(0.0770)(0.4750) = 0.8183$; bias: $-0.1 + 0.5(0.0770) = -0.0615$.
  \item $h_1$: weight from $x_1$ $0.4 \to 0.4054$, bias $0.1 \to 0.1054$; from $x_2$
        unchanged, because $x_2 = 0$.
  \item $h_2$: weight from $x_1$ $-0.3 \to -0.2923$, bias $0.2 \to 0.2077$.
\end{tight}
\textbf{Check.} A new forward pass gives $o = 0.6727$: closer to 1. The squared error
$\tfrac{1}{2}(t - o)^2$ falls from 0.0586 to 0.0536. Thousands of such small steps
are what training a network means.
\end{worked}

\section{Activation Functions and Outputs}

\dsfig{%
\begin{tikzpicture}
\begin{axis}[width=9.0cm, height=5.0cm, axis lines=middle,
  xlabel={$z$}, ylabel={activation},
  tick label style={font=\tiny}, label style={font=\scriptsize},
  xmin=-4, xmax=4, ymin=-1.2, ymax=2.2,
  legend style={font=\tiny, draw=gray!40, at={(0.02,0.97)}, anchor=north west}]
\addplot[secondary, line width=1.3pt, domain=-4:4, samples=100] {1/(1+exp(-x))};
\addlegendentry{sigmoid $\sigma(z)$: $(0, 1)$}
\addplot[purplehl, line width=1.3pt, domain=-4:4, samples=100] {tanh(x)};
\addlegendentry{tanh: $(-1, 1)$}
\addplot[accent, line width=1.3pt, domain=-4:2.2, samples=60] {max(0,x)};
\addlegendentry{ReLU $\max(0, z)$}
\end{axis}
\end{tikzpicture}}%
{Three activation functions. The sigmoid and tanh flatten out for large $|z|$, where their
gradients vanish; ReLU does not saturate for positive inputs, which is why it trains deep
networks faster.}{activations}

Mitchell's networks use sigmoids throughout. Modern practice differs in two ways.
Hidden layers use \term{ReLU}, $\max(0, z)$: its gradient is 1 for any positive
input, so it does not shrink the error signal the way a saturated sigmoid does
(the \emph{vanishing gradient} that made deep sigmoid networks hard to train). And
the output layer matches the task: a single sigmoid with cross-entropy for two
classes (exactly logistic regression, \chref{logreg}), a softmax with cross-entropy
for several, a linear unit with squared error for regression.

\section{Training Well: Overfitting and Its Cures}

A network with enough hidden units can memorise any training set, and it will if
allowed to. \dsref{early-stopping} shows the typical course: training error falls
steadily, validation error falls for a while and then rises as the network begins
fitting the noise.

\dsfig{%
\begin{tikzpicture}
\begin{axis}[width=9.6cm, height=5.0cm, axis lines=left,
  xlabel={epoch}, ylabel={cross-entropy loss},
  tick label style={font=\tiny}, label style={font=\scriptsize},
  xmin=0, xmax=150, ymin=0, ymax=1.9,
  legend style={font=\tiny, draw=gray!40, at={(0.98,0.5)}, anchor=east}]
\addplot[primary, line width=1.2pt] coordinates
  {(1,0.632) (2,0.565) (3,0.517) (4,0.482) (5,0.454) (6,0.430) (7,0.408) (8,0.387) (9,0.367) (10,0.348) (20,0.214) (30,0.110) (40,0.043) (50,0.014) (60,0.006) (70,0.003) (80,0.002) (90,0.001) (100,0.001) (110,0.001) (120,0.001) (130,0.001) (140,0.001) (150,0.001)};
\addlegendentry{training}
\addplot[accent, line width=1.2pt] coordinates
  {(1,0.669) (2,0.636) (3,0.618) (4,0.611) (5,0.610) (6,0.613) (7,0.617) (8,0.620) (9,0.621) (10,0.621) (20,0.666) (30,0.773) (40,0.922) (50,1.122) (60,1.313) (70,1.442) (80,1.537) (90,1.602) (100,1.648) (110,1.686) (120,1.718) (130,1.747) (140,1.773) (150,1.798)};
\addlegendentry{validation}
\draw[greenhl!70!black, dashed, line width=0.9pt] (axis cs:5,0) -- (axis cs:5,1.0);
\node[font=\tiny, greenhl!45!black, anchor=south west, align=left] at (axis cs:5,1.0)
     {stop here: epoch 5};
\end{axis}
\end{tikzpicture}}%
{Training and validation loss of a two-hidden-layer network on data with 20\% label noise.
Training loss goes to zero; validation loss is lowest after five epochs and rises
thereafter. Early stopping keeps the epoch-5 weights.}{early-stopping}

\rowcolors{2}{white}{lightbg}
\begin{longtable}{@{}L{3.2cm}L{6.0cm}L{5.8cm}@{}}
\toprule
\rowcolor{primary!15}
\textbf{Technique} & \textbf{What it does} & \textbf{Setting} \\
\midrule
\endhead
Early stopping & Keep the weights from the epoch with the lowest validation error &
Needs a validation set; patience of a few epochs \\
Weight decay (L2) & Add $\lambda\sum w^2$ to the loss --- ridge regression for networks
(\chref{generalization}) & \texttt{alpha} in scikit-learn; tune by CV \\
Dropout & During training, switch off each hidden unit at random with probability $p$,
so no unit can rely on another & $p \approx 0.2$--$0.5$; standard in deep learning
frameworks \\
Fewer units & A smaller network has less capacity to memorise & Tune the width by CV \\
More data, augmented & Real or synthetic variations of the training examples &
Paraphrased ticket text; small random jitter on numeric features \\
\bottomrule
\end{longtable}

\begin{conceptbox}[What the Hidden Units Learn, and What Comes Next]
Backpropagation lets the hidden layer invent its own features. Mitchell's famous
demonstration trains an 8--3--8 network to copy an 8-bit input to its output
through only 3 hidden units; the hidden units discover a 3-bit binary code for the
eight inputs, a representation nobody programmed. Deep learning is this idea
scaled up: many layers, each building features from the one below, with
architectures designed for the data --- convolutional networks for images,
recurrent networks and transformers for sequences. Those architectures are the
subject of the Deep Learning course; this chapter is their foundation.
\end{conceptbox}

\begin{pitfall}
\begin{tight}
  \item \textbf{Not scaling inputs.} Unscaled inputs saturate sigmoid units and
        make gradient descent crawl.
  \item \textbf{Initialising all weights to zero.} Every hidden unit then computes
        the same thing and receives the same update forever; initialise randomly.
  \item \textbf{Training until the training error is zero.} Watch the validation
        error and stop when it turns.
  \item \textbf{Assuming a network beats a linear model.} On small tabular data it
        often does not; in this chapter's lab it gains less than one point.
  \item \textbf{Judging from one run.} Random initialisation makes results vary;
        repeat with several seeds.
\end{tight}
\end{pitfall}

%---------------------------------------------------------------------
\begin{lab}[Backpropagation From Scratch, Then a Real Network]
\begin{lstlisting}[language=Python]
import numpy as np
from sklearn.model_selection import train_test_split
from sklearn.preprocessing import StandardScaler
from sklearn.neural_network import MLPClassifier
from sklearn.linear_model import LogisticRegression

# --- 1. XOR: no perceptron can learn it; a 2-3-1 network can -------------
X = np.array([[0, 0], [0, 1], [1, 0], [1, 1]], dtype=float)
t = np.array([[0], [1], [1], [0]], dtype=float)
sig = lambda z: 1 / (1 + np.exp(-z))
rng = np.random.default_rng(1)
W1, b1 = rng.normal(0, 1, (2, 3)), np.zeros(3)     # input -> hidden
W2, b2 = rng.normal(0, 1, (3, 1)), np.zeros(1)     # hidden -> output
eta = 1.0
for epoch in range(5001):
    h = sig(X @ W1 + b1)                           # forward pass
    o = sig(h @ W2 + b2)
    d_out = o * (1 - o) * (t - o)                  # output deltas
    d_hid = h * (1 - h) * (d_out @ W2.T)           # hidden deltas
    W2 += eta * h.T @ d_out                        # weight updates
    b2 += eta * d_out.sum(0)
    W1 += eta * X.T @ d_hid
    b1 += eta * d_hid.sum(0)
    if epoch in (0, 500, 5000):
        err = 0.5 * np.sum((t - o) ** 2)
        print(f"epoch {epoch:4d}: error {err:.4f}  "
              f"outputs {o.ravel().round(2)}")

# --- 2. staffing: too few people is late, and so is too many -------------
rng = np.random.default_rng(0)
n = 800
scope = rng.uniform(100, 1000, n)                 # story points
team = rng.uniform(2, 20, n)                      # people
Xs = np.column_stack([scope, team])
ys = (np.abs(team - scope / 60) > 5).astype(int)  # 1 = late
flip = rng.random(n) < 0.05                       # 5% recorded wrongly
ys[flip] = 1 - ys[flip]
X_tr, X_te, y_tr, y_te = train_test_split(Xs, ys, test_size=0.3,
                                          random_state=0, stratify=ys)
sc = StandardScaler().fit(X_tr)
X_tr, X_te = sc.transform(X_tr), sc.transform(X_te)
lin = LogisticRegression().fit(X_tr, y_tr)
print("logistic regression test accuracy:", round(lin.score(X_te, y_te), 3))
for alpha in (0.0001, 1.0):                        # weight decay
    net = MLPClassifier(hidden_layer_sizes=(16,), alpha=alpha,
                        max_iter=3000, random_state=0).fit(X_tr, y_tr)
    print(f"network, weight decay {alpha:g}: {net.n_iter_} epochs, "
          f"test accuracy {net.score(X_te, y_te):.3f}")
\end{lstlisting}
In part 2 a project runs late when its team is far from the size its scope needs
--- understaffed, or so overstaffed that coordination swamps the work (Brooks's
law). The late projects lie on \emph{both} sides of a diagonal band, so no single
line separates them: logistic regression manages 0.56, barely better than a coin,
while the network finds the band and scores 0.90. Now add
\texttt{early\_stopping=True} to the network and explain why accuracy
\emph{falls} so sharply.
\end{lab}

%---------------------------------------------------------------------
\begin{checkpoint}
\begin{tightnum}
  \item A perceptron has $w_0 = -1$, $w_1 = 2$, $w_2 = 1$. What does it output for
        $(0, 1)$ and $(1, 0)$?
  \item Why can a single perceptron not learn XOR, and how does one hidden layer
        fix it?
  \item In backpropagation, why does the hidden unit's delta contain the output
        unit's delta?
\end{tightnum}
\end{checkpoint}

%---------------------------------------------------------------------
\begin{chaptersummary}
\begin{tight}
  \item A perceptron thresholds a weighted sum; the perceptron rule finds a
        separating hyperplane whenever one exists.
  \item A single unit cannot represent XOR or any function that is not linearly
        separable; hidden layers can.
  \item The delta rule is gradient descent on squared error and converges even on
        non-separable data; sigmoid units make it differentiable.
  \item Backpropagation computes output deltas, passes them backwards through the
        weights to the hidden units, and updates every weight by
        $\eta\,\delta\,x$ --- the chain rule applied layer by layer.
  \item Modern networks use ReLU hidden units and task-matched outputs
        (sigmoid, softmax, linear).
  \item Control overfitting with early stopping, weight decay, dropout, smaller
        networks and more data.
\end{tight}
\end{chaptersummary}

\begin{reviewq}
  \item \textbf{(MCQ)} Which function can a single perceptron \emph{not} represent?
        (a)~AND (b)~OR (c)~NAND (d)~XOR.
  \item \textbf{(MCQ)} The derivative of the sigmoid output $o$ with respect to its
        input is: (a)~$o$ (b)~$1 - o$ (c)~$o(1 - o)$ (d)~$e^{-o}$.
  \item \textbf{(MCQ)} Early stopping chooses the weights from the epoch with the
        lowest: (a)~training error (b)~validation error (c)~test error (d)~learning
        rate.
  \item \textbf{(Short)} Contrast the perceptron training rule with the delta rule.
  \item \textbf{(Short)} Why must a network's weights be initialised randomly?
  \item \textbf{(Short)} Explain why ReLU helps train deep networks.
  \item \textbf{(Applied)} Give weights for a two-layer network of threshold units
        that computes XOR, and verify it on all four inputs.
  \item \textbf{(Applied)} Repeat the backpropagation worked example for the input
        $(0, 1)$ with target $t = 0$, starting from the original weights. Report the
        output before and after the update.
\end{reviewq}
